\documentclass[12pt, a4paper]{article}

% --- Packages ---
\usepackage[margin=2.5cm]{geometry}
\usepackage{amsmath, amssymb, amsthm}
\usepackage{booktabs}
\usepackage{array}
\usepackage{tabularx}
\usepackage{longtable}
\usepackage{hyperref}
\usepackage{xcolor}
\usepackage{titlesec}
\usepackage{enumitem}
\usepackage{microtype}
\usepackage{setspace}
\usepackage{parskip}
\usepackage{natbib}
\usepackage{caption}

% --- Hyperref setup ---
\hypersetup{
    colorlinks=true,
    linkcolor=black,
    citecolor=teal,
    urlcolor=teal,
    pdftitle={Thesis Proposal v3},
    pdfauthor={Joost}
}

% --- Theorem environments ---
\newtheorem{proposition}{Proposition}
\newtheorem{remark}{Remark}

% --- Section formatting ---
\titleformat{\section}{\large\bfseries}{}{0em}{}[\titlerule]
\titleformat{\subsection}{\normalsize\bfseries}{}{0em}{}
\titleformat{\subsubsection}{\normalsize\itshape}{}{0em}{}

% --- Column type for tables ---
\newcolumntype{L}[1]{>{\raggedright\arraybackslash}p{#1}}
\newcolumntype{C}[1]{>{\centering\arraybackslash}p{#1}}

% --- Custom colours ---
\definecolor{darkgray}{gray}{0.3}
\definecolor{medgray}{gray}{0.5}

% --- Spacing ---
\setstretch{1.15}

% ============================================================
\begin{document}
% ============================================================

% --- Title block ---
\begin{center}
    {\LARGE\bfseries Thesis Proposal v3\,---\,Final}\\[1em]
    {\large\bfseries FDR-Controlled Variable Selection under Parameter Instability:\\
    Comparing IPAD Knockoffs and OCMT with an Application\\
    to Equity Premium Prediction}\\[1.2em]
    {\normalsize BSc Econometrics \& Operations Research\,\textbar\,Erasmus University Rotterdam}\\
    {\normalsize Supervisor: Sam van Meer}\\[0.5em]
    {\color{darkgray}\small\today}
\end{center}

\vspace{0.5em}
\hrule
\vspace{1.5em}

% ============================================================
\section{Context and Motivation}
% ============================================================

Variable selection in high-dimensional regressions is a central problem in econometric forecasting. Two recent approaches address it from different angles:

\begin{enumerate}[leftmargin=*, label=\arabic*.]
    \item \textbf{IPAD} \citep{fan2020ipad}: uses Model-X knockoffs with a factor model for the covariate distribution to achieve FDR-controlled variable selection. Explicitly assumes stable parameters: \emph{``Our work has focused on the scenario of static models''} (p.\,1833).

    \item \textbf{OCMT under instability} \citep{chudik2024ocmt}: studies the One Covariate at a Time Multiple Testing procedure under time-varying parameters. Shows OCMT selects variables with non-zero \emph{average} effects and recommends full-sample selection with down-weighted forecasting. Provides no FDR control.
\end{enumerate}

Mahrad Sharifvaghefi is a co-author on both papers. These represent two lines of work by the same researcher attacking the same problem---high-dimensional variable selection for economic forecasting---from different frameworks. No paper compares them or asks: \emph{under parameter instability, when should a forecaster prefer FDR-controlled selection (IPAD) over threshold-based selection (OCMT), and can we get the best of both?}

% ============================================================
\section{Research Questions}
% ============================================================

\begin{enumerate}[leftmargin=*, label=\textbf{RQ\arabic*.}]
    \item \textbf{Diagnostic.} How does parameter instability affect the FDR control and power of IPAD knockoffs? How does it affect the false discovery behaviour of OCMT (which has no formal FDR guarantee)?

    \item \textbf{Comparative.} Under structural breaks, which method provides better variable selection for forecasting---IPAD (FDR control, assumes stability) or OCMT (handles instability, no FDR control)?

    \item \textbf{Constructive.} Can discounted IPAD---using exponentially weighted LCD statistics---bridge the gap, achieving both OCMT's robustness to breaks and IPAD's FDR control?
\end{enumerate}

% ============================================================
\section{Core Conceptual Contribution: The Dual-Null Separation}
% ============================================================

Under a structural break at $t^*$ with coefficients $\beta^{(1)}$ (pre-break) and $\beta^{(2)}$ (post-break), the relevant null hypothesis is ambiguous. Table~\ref{tab:dualnull} formalises this.

\begin{table}[h!]
\centering
\caption{Dual-null separation under a structural break}
\label{tab:dualnull}
\small
\begin{tabular}{L{2.8cm} L{5cm} L{3cm} L{3.5cm}}
\toprule
\textbf{Null} & \textbf{Definition} & \textbf{Targeted by} & \textbf{Relevant for} \\
\midrule
Local null $H_0^{(2)}$ & $\{j : \beta_j^{(2)} = 0\}$ & Discounted IPAD, Oracle & Forecasting at $T+1$ \\[4pt]
Average null $H_0^{\text{avg}}$ & $\{j : T^{-1} \sum_t \beta_{t,j} = 0\}$ & OCMT (implicitly) & Detecting ``ever-active'' variables \\[4pt]
Full-sample null $H_0^{\text{full}}$ & $\{j : \text{plim}\;\hat{\beta}_j^{\text{OLS}} = 0\}$ & Standard IPAD & Misaligned under breaks \\
\bottomrule
\end{tabular}
\end{table}

\textbf{This is not merely a technicality.} A variable active in regime 1 but null in regime 2 is a \emph{true} discovery under $H_0^{\text{avg}}$ (the OCMT perspective) but a \emph{false} discovery under $H_0^{(2)}$ (the forecasting perspective). This separation explains \emph{why} IPAD and OCMT behave differently under breaks---they answer different inferential questions. Both nulls are evaluated and reported throughout.

% ============================================================
\section{Theoretical Foundation}
% ============================================================

\subsection{Model-X Knockoffs \citep{candes2018}}

Given i.i.d.\ observations $(X_i, Y_i)$, construct knockoff variables $\tilde{X}$ satisfying:
\begin{itemize}[leftmargin=*]
    \item \emph{Pairwise exchangeability:} $(X, \tilde{X})_{\text{swap}(S)} \stackrel{d}{=} (X, \tilde{X})$ for all $S \subseteq \{1,\ldots,p\}$
    \item \emph{Conditional independence:} $\tilde{X} \perp\!\!\!\perp Y \mid X$
\end{itemize}

Feature importance statistics $W_j$ satisfy the sign-flip property, and the knockoff+ filter at target level $q$ provides finite-sample FDR control: $\text{FDR} \leq q$.

\textbf{Key point.} The knockoff construction depends only on $P_X$, not on $P_{Y|X}$. Therefore, if $P_X$ is stable but $\beta_t$ changes, knockoffs remain valid---only the signal detection step is affected.

\subsection{IPAD \citep{fan2020ipad}}

Models covariates via a latent factor model $X_t = BF_t + U_t$. Estimates $\Sigma_X = BB' + \Sigma_U$ via PCA (with Ahn--Horenstein eigenvalue ratio for $\hat{r}$), enabling approximate knockoff construction. Uses the LCD statistic:
\[
    W_j = |\hat{\beta}_j| - |\hat{\beta}_{j+p}|,
\]
where $(\hat{\beta}_1, \ldots, \hat{\beta}_{2p})$ is the Lasso estimate on the augmented design $[X, \tilde{X}]$. Provides asymptotic FDR control. \textbf{Assumes stable $\beta$.}

\subsection{OCMT \citep{chudik2018, chudik2024ocmt}}

Tests each covariate individually: regress $Y$ on $X_j$, obtain $t$-statistic, threshold at
\[
    c_{p,T} = \Phi^{-1}\!\left(1 - \frac{p}{2\delta_T}\right), \quad p/\delta_T \to 0.
\]
Under stability, OCMT consistently selects all signals and no noise. Under instability, OCMT selects variables with non-zero \emph{average} marginal effects. \textbf{No FDR control is provided or claimed.}

% ============================================================
\section{The Gap}
% ============================================================

\begin{table}[h!]
\centering
\caption{Literature gap: what is known and what remains open}
\label{tab:gap}
\small
\begin{tabular}{L{3.5cm} L{4.8cm} L{5cm}}
\toprule
\textbf{Source} & \textbf{What it says} & \textbf{What remains open} \\
\midrule
\citet{fan2020ipad}, p.\,1833 & ``Our work has focused on the scenario of static models'' & What happens to IPAD FDR when $\beta$ changes? \\[4pt]
\citet{chudik2024ocmt}, \S5 & OCMT selects variables with non-zero average effects; recommends down-weighting at forecasting stage only & Can selection itself be localised? What is OCMT's empirical FDR? \\[4pt]
Chi et al.\ (2025) --- TSKI & Handles serial dependence; assumes stable $\beta$ & Orthogonal to this thesis (dependence $\neq$ instability) \\[4pt]
Liu, Sun \& Ke (2024) --- GKnockoff & Uses knockoffs to \emph{detect} change points & Does not do variable selection \emph{despite} breaks \\[4pt]
\citet{barber2020} & Robust knockoffs: FDR $\leq q \cdot e^\epsilon$ when $P_X$ is misspecified & Does not address $\beta$ instability (different misspecification) \\
\bottomrule
\end{tabular}
\end{table}

\textbf{No paper compares IPAD and OCMT under instability. No paper studies discounted knockoff statistics. No paper formalises the dual-null issue.}

% ============================================================
\section{Methods}
% ============================================================

\subsection{Method 1: Standard IPAD (Baseline)}
PCA-based factor estimation with Ahn--Horenstein (2013) eigenvalue ratio for $\hat{r}$; second-order approximate knockoff construction from $\hat{\Sigma}_X$; LCD statistics via Lasso (cross-validated $\lambda$) on augmented design $[X, \tilde{X}]$; knockoff+ filter at $q = 0.20$.

\subsection{Method 2: OCMT (Baseline)}
Univariate OLS for each predictor; threshold at $c_{p,T} = \Phi^{-1}(1 - p/(2T^{1.1}))$; empirical FDR reported (informative even without formal guarantee).

\subsection{Method 3: Discounted IPAD (Proposed Modification)}

\textbf{Idea.} Replace the standard Lasso with an exponentially weighted Lasso:
\[
    \hat{\beta}^{(\lambda,\delta)} = \underset{b}{\arg\min}\ \sum_{t=1}^T \delta^{T-t}\bigl(Y_t - [X_t, \tilde{X}_t]'\, b\bigr)^2 + \lambda\|b\|_1.
\]

\textbf{Implementation.} Multiply observation $t$ by weight $\sqrt{\delta^{T-1-t}}$, then fit standard Lasso on the reweighted data.

\begin{proposition}[Sign-flip property of weighted LCD]
If $w_t > 0$ are deterministic weights constructed independently of $\tilde{X}$, then the weighted LCD statistic $W_j^{(w)} = |\hat{\beta}_j^{(w)}| - |\hat{\beta}_{j+p}^{(w)}|$ satisfies the sign-flip property.
\end{proposition}

\begin{proof}
Define reweighted data $(Y^w, Z^w)$ where $Y_t^w = w_t Y_t$ and $Z_t^w = w_t[X_t, \tilde{X}_t]$. Since $w_t$ is deterministic and independent of $\tilde{X}$:
\begin{enumerate}[label=(\roman*)]
    \item \emph{Conditional independence:} $\tilde{X} \perp\!\!\!\perp Y \mid X$ implies $Z^w_{\cdot,j+p} \perp\!\!\!\perp Y^w \mid Z^w_{\cdot,1:p}$ (weighting preserves conditional independence).
    \item \emph{Column exchangeability:} $(X,\tilde{X})_{\text{swap}(S)} \stackrel{d}{=} (X,\tilde{X})$ implies $(Z^w_{\cdot,1:p}, Z^w_{\cdot,p+1:2p})_{\text{swap}(S)} \stackrel{d}{=} (Z^w_{\cdot,1:p}, Z^w_{\cdot,p+1:2p})$ since the same row-scaling applies to both original and knockoff columns.
    \item Therefore $W_j^{(w)}(Y^w,[Z^w]_{\text{swap}(S)}) = -W_j^{(w)}(Y^w,Z^w)$ for $j \in S$, by the standard sign-flip argument. \qedhere
\end{enumerate}
\end{proof}

\textbf{Effective sample size.} The weighted Lasso behaves as if operating on
\[
    n_{\text{eff}}(\delta) = \frac{1 - \delta^{2T}}{1 - \delta^2}
\]
effective observations. For $\delta = 0.97$, $T = 500$: $n_{\text{eff}} \approx 17$. This analytically predicts the power--localisation trade-off.

\textbf{Discount factor selection.} Grid $\delta \in \{0.95, 0.97, 0.98, 0.99, 1.00\}$. In the empirical application, select $\delta$ by minimising out-of-sample forecast MSFE over a 60-month validation window.

\subsection{Method 4: Rolling-Window IPAD}
Knockoff construction from full sample ($P_X$ stable); LCD statistics computed within a trailing window of size $h \in \{0.3T, 0.5T, 0.7T\}$; same knockoff+ filter applied to windowed LCD values.

\subsection{Method 5: Down-Weighted OCMT \citep{chudik2024ocmt}}
Full-sample OCMT for variable selection (unweighted); exponentially weighted OLS for post-selection forecasting. This is the approach recommended in the source paper.

\subsection{Method 6: Oracle Benchmark}
Knows the break date $t^*$; uses full sample for knockoff construction ($P_X$ stable); restricts LCD statistics to post-break observations only. Provides an upper bound on what any localised method can achieve.

% ============================================================
\section{Data Generating Processes}
% ============================================================

\textbf{Covariate model (stable throughout):}
\[
    X_t = BF_t + U_t, \quad F_t \sim N(0, I_r),\quad U_t \sim N(0, \operatorname{diag}(\sigma_1^2, \ldots, \sigma_p^2))
\]
with $r = 3$ factors, loadings $B_{jk} \sim U(-1,1)$, idiosyncratic variances $\sigma_j^2 \sim U(0.5, 1.5)$.

\textbf{Response model:}
\[
    Y_t = X_t'\beta_t + \varepsilon_t, \quad \varepsilon_t \sim N(0,1).
\]

\begin{table}[h!]
\centering
\caption{Simulation scenarios (prioritised)}
\label{tab:scenarios}
\small
\begin{tabular}{C{1.5cm} L{3cm} L{4cm} L{5cm}}
\toprule
\textbf{Priority} & \textbf{Scenario} & \textbf{Mechanism} & \textbf{What it tests} \\
\midrule
Primary & Signal entry/exit & $S^{(1)} \neq S^{(2)}$: some variables active only pre-break, others only post-break & FDR inflation under local null; dual-null divergence \\[4pt]
Primary & Stable (baseline) & $\beta$ constant & Sanity check: verify both methods work as advertised \\[4pt]
Secondary & Magnitude shift & $\beta^{(2)} = c\cdot\beta^{(1)}$, same support & Power under signal dilution \\[4pt]
Secondary & Sign reversal & $\beta^{(2)} = -\beta^{(1)}$ & Average effect $= 0$; OCMT should fail \\
\bottomrule
\end{tabular}
\end{table}

\textbf{Signal entry/exit specification.} $|S^{(1)}| = |S^{(2)}| = s$; overlap $|S^{(1)} \cap S^{(2)}| = s/2$; non-zero coefficients $\beta_j = A/\sqrt{s}$ where $A$ is calibrated to target SNR.

% ============================================================
\section{Simulation Design}
% ============================================================

\begin{table}[h!]
\centering
\caption{Dimensionality configurations}
\label{tab:dimconfigs}
\small
\begin{tabular}{C{2cm} C{3cm} L{6cm}}
\toprule
\textbf{Config} & $\boldsymbol{(T,p,s)}$ & \textbf{Role} \\
\midrule
Low-dim & $(500, 50, 5)$ & Sanity check; fast iteration \\
\textbf{Baseline} & $\boldsymbol{(500, 200, 20)}$ & \textbf{Primary results} \\
High-dim & $(300, 300, 15)$ & $p = T$ boundary \\
\bottomrule
\end{tabular}
\end{table}

\textbf{Primary grid:}
\begin{itemize}[leftmargin=*]
    \item Break location: $\tau \in \{0.5, 0.7\}$
    \item Break magnitude (overlap fraction): $\in \{1.0, 0.75, 0.5, 0.25, 0.0\}$
    \item SNR: $\in \{2, 4\}$
    \item Discount factor: $\delta \in \{0.95, 0.97, 0.99, 1.00\}$
    \item Replications: $M = 500$
\end{itemize}

Total primary configurations: $2 \times 5 \times 2 \times 4 = 80$ cells $\times$ 6 methods $\times$ 500 reps. At approximately 0.5s per replication (Lasso on $n=500, p=400$), this is $\approx$67 CPU-hours, parallelisable to under 8 hours on 8 cores.

\textbf{Evaluation metrics:}

\emph{Primary (reported for all methods):}
\begin{enumerate}[leftmargin=*, label=(\arabic*)]
    \item Empirical FDR against local null $H_0^{(2)}$ (the forecasting-relevant null)
    \item Empirical FDR against average null $H_0^{\text{avg}}$
    \item Power (TPR) against local null $H_0^{(2)}$
    \item Out-of-sample $R^2_{\text{OOS}}$ (one-step-ahead forecast using post-selection OLS)
\end{enumerate}

\emph{Secondary:}
\begin{enumerate}[leftmargin=*, label=(\arabic*), start=5]
    \item Selection stability: Jaccard similarity $J(\hat{S}_i, \hat{S}_j)$ averaged over replication pairs
    \item FDR conditional on discovery: $\text{FDR} \mid |\hat{S}| > 0$ (avoids the $0/0$ issue)
\end{enumerate}

\textbf{Key results sections:}

\begin{description}[leftmargin=2em]
    \item[\S4.1 Baseline replication.] Under stability: verify IPAD controls FDR at $q = 0.20$; verify OCMT's empirical FDR.
    \item[\S4.2 Breakdown analysis.] Trace the FDR of standard IPAD as break severity increases. Plot FDR vs.\ overlap fraction, stratified by SNR.
    \item[\S4.3 Head-to-head comparison.] IPAD vs.\ OCMT: FDR--power frontiers under each scenario.
    \item[\S4.4 Discounted IPAD.] Show $\delta < 1$ restores FDR control at the cost of power; plot breakdown threshold curve in (overlap, $\delta$) space; compare empirical power to the $n_{\text{eff}}$ prediction.
    \item[\S4.5 The dual null in action.] Show that the same run of discounted IPAD has different empirical FDR depending on whether evaluated against $H_0^{(2)}$ or $H_0^{\text{avg}}$.
    \item[\S4.6 Forecasting comparison.] $R^2_{\text{OOS}}$ across methods; DM tests for pairwise forecast comparisons.
\end{description}

% ============================================================
\section{Empirical Application: Equity Premium Prediction}
% ============================================================

\textbf{Data.} Monthly equity premium $r_{t+1}^e = r_{t+1}^{\text{mkt}} - r_f$; 14 standard \citet{welch2008} variables plus 3 lags of each ($p = 56$); sample 1950:01--2023:12; OOS evaluation 1966:01--2023:12 (expanding window, minimum 15 years in-sample).

\textbf{Design.} At each forecast origin $t$, apply: (i) IPAD, (ii) OCMT, (iii) Discounted IPAD ($\delta$ chosen by validation MSFE over prior 60 months). Post-selection forecast: OLS on selected variables (or weighted OLS for discounted methods).

\textbf{Outputs.}
\begin{enumerate}[leftmargin=*, label=(\arabic*)]
    \item \emph{Selection heatmap:} variables (rows) $\times$ time (columns), coloured by selection frequency across 100 knockoff repetitions.
    \item \emph{OOS $R^2$ comparison:} cumulative $R^2_{\text{OOS}}$ for each method vs.\ historical mean benchmark.
    \item \emph{Break narrative:} do variable selections shift around known instability episodes (1973 oil shock, 2008 GFC, 2020 COVID)?
\end{enumerate}

% ============================================================
\section{Maintained Assumptions}
% ============================================================

\begin{enumerate}[leftmargin=*, label=\textbf{A\arabic*.}]
    \item \textbf{Stable $P_X$.} The covariate distribution does not change at the break. Only $\beta_t$ changes. This isolates the coefficient instability channel and ensures knockoff validity.
    \item \textbf{Single break.} The simplest form of instability. Discussion covers extensions to multiple breaks and gradual drift.
    \item \textbf{Cross-sectional independence across $t$.} Observations are i.i.d.\ draws from the factor model within each period. Serial dependence is handled by TSKI (explicitly out of scope).
    \item \textbf{Known target FDR level.} $q = 0.20$ throughout (standard in the knockoff literature).
\end{enumerate}

% ============================================================
\section{Expected Contributions}
% ============================================================

\begin{table}[h!]
\centering
\caption{Summary of expected contributions}
\label{tab:contributions}
\small
\begin{tabular}{C{0.4cm} L{5.5cm} C{2cm} L{4.5cm}}
\toprule
\textbf{\#} & \textbf{Contribution} & \textbf{Type} & \textbf{Strength} \\
\midrule
1 & First comparison of IPAD and OCMT under parameter instability & MC Empirical & Main result: who wins when? \\[4pt]
2 & Dual-null separation: IPAD and OCMT target different nulls & Conceptual & Framework-level insight; claimed regardless of MC outcomes \\[4pt]
3 & Discounted LCD statistics with formal sign-flip verification & Methodological & Simple but novel; extends knockoff toolkit \\[4pt]
4 & Breakdown threshold characterisation: when does IPAD's FDR guarantee fail? & Diagnostic & Practical guidance for applied researchers \\[4pt]
5 & Equity premium application with time-varying selections & Empirical & Connects to macro-finance literature \\
\bottomrule
\end{tabular}
\end{table}

% ============================================================
\section{Timeline}
% ============================================================

\begin{table}[h!]
\centering
\caption{8-week thesis timeline}
\label{tab:timeline}
\small
\begin{tabular}{C{1.2cm} L{5cm} L{7cm}}
\toprule
\textbf{Week} & \textbf{Milestone} & \textbf{Deliverable} \\
\midrule
1 & Finalise proposal. Verify IPAD baseline. Implement OCMT. Run stable-DGP sanity check. & Working comparison pipeline; Table 1 (baseline FDR/power). \\[3pt]
2 & Implement break DGPs. Run IPAD under instability. Document FDR inflation. & Figure: FDR vs.\ overlap fraction. \S4.1--4.2 draft results. \\[3pt]
3 & Implement discounted IPAD + rolling-window IPAD. Run against oracle. & \S4.3--4.4 results. Breakdown threshold figure. \\[3pt]
4 & Full simulation grid (primary). Generate all tables/figures. & Complete simulation results. \\[3pt]
5 & Secondary grid (high-dim, sign reversal). Begin empirical application. & Appendix tables. Application data pipeline working. \\[3pt]
6 & Finish empirical application. Write \S4 (Results) and \S5 (Application). & Draft results + application chapters. \\[3pt]
7 & Write \S1 (Intro), \S2 (Literature), \S3 (Methods), \S6 (Conclusion). & Full draft. \\[3pt]
8 & Revisions, polish, buffer for supervisor feedback. & Final submission. \\
\bottomrule
\end{tabular}
\end{table}

% ============================================================
\section{Scope Management}
% ============================================================

\textbf{Hard scope (must deliver):}
\begin{itemize}[leftmargin=*]
    \item Simulation comparison: IPAD vs.\ OCMT vs.\ Discounted IPAD vs.\ Oracle
    \item Primary grid: baseline + signal entry/exit scenario, $(500, 200, 20)$
    \item Sign-flip proof for discounted LCD
    \item Dual-null evaluation
    \item Empirical application (streamlined: 1 expanding window, 3 methods)
\end{itemize}

\textbf{Soft scope (deliver if time allows):}
\begin{itemize}[leftmargin=*]
    \item High-dimensional configuration; rolling-window IPAD variant
    \item Magnitude shift and sign reversal scenarios
    \item Derandomisation (run each method 50$\times$ with different knockoff seeds; report stability)
    \item DM tests for forecast comparison
\end{itemize}

\textbf{Explicitly out of scope:}
\begin{itemize}[leftmargin=*]
    \item Formal asymptotic theory for discounted IPAD FDR
    \item Serial dependence in $X_t$ (handled by TSKI)
    \item Joint $P_X$ and $\beta$ instability; deep generative knockoff construction
    \item Multiple simultaneous breaks (discussed qualitatively only)
\end{itemize}

% ============================================================
\section{Risk Mitigation}
% ============================================================

\begin{table}[h!]
\centering
\caption{Risk assessment and mitigation}
\label{tab:risks}
\small
\begin{tabular}{L{4cm} C{2cm} L{7.5cm}}
\toprule
\textbf{Risk} & \textbf{Likelihood} & \textbf{Mitigation} \\
\midrule
Discounted IPAD does not control FDR & Medium & The comparison is the main contribution. A negative result---``discounting helps power but does not fully restore FDR''---is informative. Report oracle as upper bound. \\[4pt]
Computation too slow & Low & Primary grid is $\approx$67 CPU-hours. Parallelise on 8 cores ($<$8 hrs). Cut reps to $M=200$ if needed. \\[4pt]
OCMT's FDR is accidentally good & Low & Would be an interesting finding. Report honestly; attribute to Bonferroni-like conservatism in the critical value. \\[4pt]
Empirical application shows no instability & Medium & Simulation study stands alone. ``Methods agree $\Rightarrow$ stability era'' is itself a finding. \\[4pt]
Supervisor wants formal theory & Medium & Offer the $n_{\text{eff}}$ analytical prediction plus sign-flip proof as the theoretical component. Frame as ``Monte Carlo evidence motivating future theoretical development.'' \\
\bottomrule
\end{tabular}
\end{table}

% ============================================================
\section{What Makes This a Strong Thesis}
% ============================================================

\begin{enumerate}[leftmargin=*, label=\arabic*.]
    \item \textbf{Dual contribution structure.} Even if discounted IPAD ``fails,'' the comparison and dual-null insight remain. The thesis has a floor.

    \item \textbf{Explicit connection between two published papers by the same co-author.} Examiners immediately see why this matters.

    \item \textbf{Clean separation of concerns.} $P_X$ instability ($\to$ TSKI) vs.\ $\beta$ instability ($\to$ this thesis). The student can explain exactly what is and is not in scope.

    \item \textbf{Analytical prediction meets simulation.} The $n_{\text{eff}}(\delta)$ formula gives a testable quantitative prediction---if the simulation matches, it validates the mechanism; if it diverges, it reveals break contamination.

    \item \textbf{Real data grounds the findings.} The equity premium application connects the simulation study to a well-known instability narrative in macro-finance.

    \item \textbf{Reproducibility.} Python code, fixed seeds, specified grid. Any examiner can replicate.
\end{enumerate}

% ============================================================
\section*{References}
% ============================================================

\begingroup
\small
\bibliographystyle{apalike}

% Inline bibliography (no .bib file required to compile)
\begin{thebibliography}{99}

\bibitem[Ahn and Horenstein(2013)]{ahn2013}
Ahn, S.~C. and Horenstein, A.~R. (2013).
\newblock Eigenvalue ratio test for the number of factors.
\newblock \emph{Econometrica}, 81(3):1203--1227.

\bibitem[Barber and Cand\`{e}s(2015)]{barber2015}
Barber, R.~F. and Cand\`{e}s, E.~J. (2015).
\newblock Controlling the false discovery rate via knockoffs.
\newblock \emph{The Annals of Statistics}, 43(5):2055--2085.

\bibitem[Barber et al.(2020)]{barber2020}
Barber, R.~F., Cand\`{e}s, E.~J., and Samworth, R.~J. (2020).
\newblock Robust inference with knockoffs.
\newblock \emph{The Annals of Statistics}, 48(3):1409--1431.

\bibitem[Cand\`{e}s et al.(2018)]{candes2018}
Cand\`{e}s, E., Fan, Y., Janson, L., and Lv, J. (2018).
\newblock Panning for gold: model-$\mathcal{X}$ knockoffs for high dimensional controlled variable selection.
\newblock \emph{Journal of the Royal Statistical Society: Series B}, 80(3):551--577.

\bibitem[Chi et al.(2025)]{chi2025}
Chi, C.~M., Fan, Y., Ing, C.~K., and Lv, J. (2025).
\newblock Tera-scale knockoff inference via aggregation.
\newblock \emph{arXiv:2112.09851}.

\bibitem[Chudik et al.(2018)]{chudik2018}
Chudik, A., Pesaran, M.~H., and Sharifvaghefi, M. (2018).
\newblock One covariate at a time, multiple testing and growth.
\newblock \emph{Journal of Econometrics}, 203(2):292--310.

\bibitem[Chudik et al.(2024)]{chudik2024ocmt}
Chudik, A., Pesaran, M.~H., and Sharifvaghefi, M. (2024).
\newblock Variable selection and forecasting in high dimensional linear regressions with parameter instability.
\newblock \emph{Journal of Econometrics}.

\bibitem[Fan et al.(2020)]{fan2020ipad}
Fan, Y., Lv, J., Sharifvaghefi, M., and Uematsu, Y. (2020).
\newblock IPAD: Stable interpretable forecasting with knockoffs inference.
\newblock \emph{Journal of the American Statistical Association}, 115(532):1822--1834.

\bibitem[Ke et al.(2024)]{ke2024}
Ke, Z.~T., Liu, X., and Ma, Y. (2024).
\newblock Power of knockoff: the impact of ranking algorithm, augmented design, and symmetric statistic.
\newblock \emph{Journal of Machine Learning Research}, 25(3).

\bibitem[Liu et al.(2024)]{liu2024}
Liu, X., Sun, Q., and Ke, Z.~T. (2024).
\newblock GKnockoff: testing hypotheses with controlled FDR in the presence of structural breaks.
\newblock \emph{Journal of Econometrics}, 239(2).

\bibitem[Nguyen et al.(2020)]{nguyen2020}
Nguyen, T.~B., Chevalier, J.~A., Thirion, B., and Arlot, S. (2020).
\newblock Aggregation of multiple knockoffs.
\newblock In \emph{Proceedings of the 37th International Conference on Machine Learning (ICML)}.

\bibitem[Ren and Barber(2023)]{ren2023}
Ren, Z. and Barber, R.~F. (2023).
\newblock Derandomised knockoffs: leveraging e-values for false discovery rate control.
\newblock \emph{Journal of the American Statistical Association}, 118(542).

\bibitem[Romano et al.(2020)]{romano2020}
Romano, Y., Sesia, M., and Cand\`{e}s, E.~J. (2020).
\newblock Deep knockoffs.
\newblock \emph{Journal of the American Statistical Association}, 115(532):1861--1872.

\bibitem[Rossi(2021)]{rossi2021}
Rossi, B. (2021).
\newblock Forecasting in the presence of instabilities: how we know whether models predict well and how to improve them.
\newblock \emph{Journal of Economic Literature}, 59(4):1135--1190.

\bibitem[Stock and Watson(2002)]{stockwatson2002}
Stock, J.~H. and Watson, M.~W. (2002).
\newblock Forecasting using principal components from a large number of predictors.
\newblock \emph{Journal of the American Statistical Association}, 97(460):1167--1179.

\bibitem[Wei et al.(2025)]{wei2025}
Wei, Y., Ke, Z.~T., and Zhang, C.~H. (2025).
\newblock FAR-knockoff: a powerful method for false discovery rate control with knockoffs.
\newblock Manuscript.

\bibitem[Welch and Goyal(2008)]{welch2008}
Welch, I. and Goyal, A. (2008).
\newblock A comprehensive look at the empirical performance of equity premium prediction.
\newblock \emph{The Review of Financial Studies}, 21(4):1455--1508.

\end{thebibliography}
\endgroup

\end{document}